% !PRINTED 280
\weberpage{280}
%<<CONT
in which the $q_1, q_2 \ldots q_{m-1}$, and likewise the $f_1, f_2 \ldots f_m$, are integral rational functions of $x$; the degrees of the functions $f, f_1, f_2, \ldots f_m$ \gesperrt{decrease}, and one can therefore continue the operation until $f_m$ is constant, or at least no longer vanishes in the interval under consideration. That $f_m$ cannot become zero is a consequence of the assumption that $f$ and $f_1$ are without common divisor.

That in the sequence
\weq{3}{f,\, f_1,\, f_2 \ldots f_m}
one then really has a Sturm chain follows immediately on going through the criteria of \S.\,86, 1. to 4. For if two consecutive ones among the functions (3) vanish simultaneously, then by (2) all the following ones vanish as well; but this is impossible, because $f_m$ is different from zero.

If however $f_\nu(x) = 0$, then from (2) it follows that
\[ f_{\nu-1}(x) = -f_{\nu+1}(x), \]
with which all the requirements of \S.\,86 are satisfied.

It is, as goes without saying, permitted to multiply the functions of the sequence (3) by positive, e.g. constant, factors, without their ceasing to form a \gesperrt{Sturm} chain.

For $n=2$ we can accordingly take as the Sturm functions the following:
\[
\begin{gathered}
f(x) = x^2 + a x + b\\
f_1(x) = 2x + a\\
f_2(x) = a^2 - 4b.
\end{gathered}
\]

\begin{center}
\S.\,90.

\gesperrt{Hermite's Solution of Sturm's Problem.}
\end{center}

Another way of solving Sturm's problem, which leads to simpler results, at least as regards the carrying out of the calculation in detail, has been taken by \gesperrt{Hermite}, who connects it with the law of inertia of quadratic forms and the Tschirnhausen transformation\footnote{\gesperrt{Hermite}, Remarques sur le th\'eor\`eme de M. Sturm. Comptes rendus of the Paris Academy, T.\,36 (1853).}.
